\section{为什么先读头文件}

头文件是系统对外承诺的最小合同。实现可以重写，kernel 可以替换，benchmark 可以补充平台路径，但头文件里暴露的数据结构、函数入口和返回字段决定了测试、CLI、工具脚本和章节讲解怎样调用这套工程。

\begin{keyidea}
读 LCQI 时不要从最长的 \code{gpt2\_reference.cpp} 开始。先读公共头文件，明确“模型如何表示”“推理返回什么”“trace 记录什么”“KV cache 怎样寻址”“tokenizer 的坏输入怎样处理”，再进入实现。这样读到长循环和 shape 检查时，知道它们在维护哪份合同。
\end{keyidea}

\section{Tiny MLP 与推理入口}

\filepath{model.hpp} 定义 tiny MLP 的模型对象、权重、bias、scale 和加载入口。它的责任是把文本模型资产解释成 C++ 对象，不负责执行推理。

\lstinputlisting[language=C++,caption={include/lcqi/model.hpp}]{../../../cpu-volume-3/labs/linux_cpu_inference/include/lcqi/model.hpp}

\filepath{inference.hpp} 定义推理结果和 \code{run\_inference} 入口。返回值包含 hidden、logits 和 predicted class，测试可以用这些字段定位是加载错、kernel 错，还是 argmax 错。

\lstinputlisting[language=C++,caption={include/lcqi/inference.hpp}]{../../../cpu-volume-3/labs/linux_cpu_inference/include/lcqi/inference.hpp}

\section{Int8 Kernel、Tokenizer 与 Safetensors}

int8 kernel 头文件定义 packed linear、deterministic fixture、scalar/workspace/AVX2 路径和误差比较。注意 AVX2 是运行时能力，不是源码存在就等于当前机器已验证。

\lstinputlisting[language=C++,caption={include/lcqi/int8\_kernels.hpp}]{../../../cpu-volume-3/labs/linux_cpu_inference/include/lcqi/int8_kernels.hpp}

F32 kernel 头文件定义 GPT-2 热路径使用的点积和 row-range 边界。读这里要区分四层：\code{dot\_f32()} 是 checked API，先确认两个 span 长度一致；\code{dot\_f32\_unchecked()} 是单行点积热路径入口；\code{linear\_f32\_rows\_unchecked()} 一次计算一段输出行，用在 QKV、attention projection 和 MLP；\code{max\_dot\_f32\_rows\_unchecked()} 一次扫描一段 vocab 行并返回局部最大值，用在 logits-free greedy \code{lm\_head}。AVX2/FMA 是否可用是运行时能力，不是 CMake 把文件编进来就等于一定执行。

\lstinputlisting[language=C++,caption={include/lcqi/f32\_kernels.hpp}]{../../../cpu-volume-3/labs/linux_cpu_inference/include/lcqi/f32_kernels.hpp}

tokenizer 头文件固定 BOS/EOS/UNK、词表和 template hash。prompt 进入模型前，必须先变成稳定 token id；未知 token 行为也必须被写进合同。

\lstinputlisting[language=C++,caption={include/lcqi/tokenizer.hpp}]{../../../cpu-volume-3/labs/linux_cpu_inference/include/lcqi/tokenizer.hpp}

safetensors 头文件固定 dtype、shape、offset 和 payload 边界。它负责模型字节怎样被解释，不应该和数学推理逻辑混在一起。

\lstinputlisting[language=C++,caption={include/lcqi/safetensors.hpp}]{../../../cpu-volume-3/labs/linux_cpu_inference/include/lcqi/safetensors.hpp}

\section{GGUF 与 Q4\_K 合同}

GGUF 头文件固定 llama.cpp 生态模型文件的读取边界：版本、metadata、alignment、tensor 名、shape、ggml type 和 tensor payload offset。它也保留 tokenizer 需要的字符串数组 metadata，例如 \code{tokenizer.ggml.tokens} 和 \code{tokenizer.ggml.merges}。它不负责模型图执行，只把“真实 GGUF 文件里有哪些 tensor、每个 tensor 在哪里、哪些 metadata 能支撑 tokenizer”变成可检查对象。

\lstinputlisting[language=C++,caption={include/lcqi/gguf.hpp}]{../../../cpu-volume-3/labs/linux_cpu_inference/include/lcqi/gguf.hpp}

GGML tensor 头文件是 GGUF 与模型数学之间的格式层。SmolLM2 的 \code{Q4\_K\_M} 文件内部并不是所有矩阵都是 \code{Q4\_K}：远端 caw 扫描显示它包含 \code{F32}、\code{Q8\_0}、\code{Q5\_0}、\code{Q6\_K} 和 \code{Q4\_K}。因此完整模型加载不能只写一个 Q4\_K 解码器，必须把这些格式统一变成后续执行层能消费的张量。

\lstinputlisting[language=C++,caption={include/lcqi/ggml\_tensors.hpp}]{../../../cpu-volume-3/labs/linux_cpu_inference/include/lcqi/ggml_tensors.hpp}

Q4\_K 头文件固定本轮量化切片的数学合同。一个 \code{Q4\_K} block 解码出 256 个权重，\code{Q8\_K} input 是每 256 个激活临时量化一次。这里暴露的 \code{matvec\_q4\_k\_q8} 是单算子入口，不是完整 LLaMA-like decoder；完整 decoder 还要接 RMSNorm、RoPE、attention、SwiGLU、KV cache 和 sampler。

\lstinputlisting[language=C++,caption={include/lcqi/q4\_k.hpp}]{../../../cpu-volume-3/labs/linux_cpu_inference/include/lcqi/q4_k.hpp}

LLaMA GGUF 头文件把真实 GGUF 文件接到 reference decoder 数学层。当前实现明确叫 \code{f32\_dequantized\_reference}：加载时把量化权重解到 f32，先建立端到端 correctness baseline，再逐步把热点替换为量化直算 micro-kernel。这个边界很重要，因为它防止把“能跑真实模型”误写成“已经拥有 llama.cpp 同等级的低比特执行引擎”。

\lstinputlisting[language=C++,caption={include/lcqi/llama\_gguf.hpp}]{../../../cpu-volume-3/labs/linux_cpu_inference/include/lcqi/llama_gguf.hpp}

\section{Reference Decoder 与 GPT-2 合同}

reference decoder 头文件定义张量、层配置、trace checkpoint、采样结果和 decode 函数。它是定位推理错误的主要合同，不只返回最终 token。

\lstinputlisting[language=C++,caption={include/lcqi/reference\_decoder.hpp}]{../../../cpu-volume-3/labs/linux_cpu_inference/include/lcqi/reference_decoder.hpp}

GPT-2 reference 头文件定义 GPT-2 配置、byte-level BPE tokenizer、HuggingFace safetensors 加载、forward、cached forward 和 greedy generation。它单独存在，是因为 GPT-2 使用 LayerNorm、绝对位置嵌入、packed QKV、GELU 和 tied lm head，不能直接塞进 LLaMA-like reference decoder。读 \code{Gpt2CachedGreedyDecoder} 时要区分三个入口：\code{step()} 推进一个 token 并返回 greedy 预测；\code{step\_with\_logits()} 保留完整 logits 给测试和对照；\code{advance\_without\_prediction()} 只推进 embedding、Transformer layers 和 KV cache，用在 prompt prefill 的前缀 token 上。第三个入口不是“少算一点但可能改变结果”，而是删掉不会被任何 caller 使用的预测结果。

\lstinputlisting[language=C++,caption={include/lcqi/gpt2\_reference.hpp}]{../../../cpu-volume-3/labs/linux_cpu_inference/include/lcqi/gpt2_reference.hpp}
